\documentclass[12pt]{article}
\usepackage[utf8]{inputenc}
\usepackage[portuguese]{babel}
\usepackage{amsmath}
\usepackage{graphicx}
\usepackage{geometry}

\geometry{a4paper, margin=1in}

\begin{document}

% ──────────────────────────────────────────────
\section*{Racional da Modelagem}

O sistema é representado por um conjunto de estados mutuamente exclusivos.
Em cada instante $t$, o sistema ocupa exatamente um desses estados, e
associa-se a cada estado $i$ uma probabilidade $P_i(t)$, que expressa a
chance de o sistema encontrar-se naquele estado no instante $t$. Por
construção, essas probabilidades satisfazem a condição de normalização:
\[
    \sum_{i} P_i(t) = 1, \quad \forall\, t \geq 0.
\]

A evolução de $P_i(t)$ é determinada pelas taxas de transição entre os
estados. Para um intervalo de tempo suficientemente pequeno $\Delta t$, a
probabilidade de ocorrer mais de uma transição é da ordem de
$O(\Delta t^2)$ e, portanto, desprezível. Com isso, a probabilidade de o
sistema estar em um estado no instante $t + \Delta t$ é a soma das
contribuições de cada estado em $t$, ponderada pela respectiva
probabilidade de transição em $\Delta t$.

% ──────────────────────────────────────────────
\newpage
\section{Definição dos Estados}

\begin{itemize}
    \item \textbf{$M_1$}: Operação plena com redundância ativa (dois módulos funcionando).
    \item \textbf{$M_2$}: Operação degradada com reparo corretivo em andamento (falha de um módulo detectada com cobertura $C$).
    \item \textbf{$FC$}: Falha coberta aguardando reparo corretivo (sistema ainda operacional, um módulo em manutenção).
    \item \textbf{$NC$}: Falha não coberta aguardando reparo preventivo (sistema ainda operacional, reparo mais lento).
    \item \textbf{$F$}: Falha total do sistema (indisponível).
\end{itemize}

% ──────────────────────────────────────────────
\newpage
\section{Equações de Transição em $\Delta t$}

Para cada estado, listam-se as probabilidades de transição durante um
intervalo $\Delta t$ suficientemente pequeno.

\subsection*{Estado $M_1$}
\begin{itemize}
    \item $P(M_1 \to M_2)   = \lambda C \,\Delta t$
    \item $P(M_1 \to FC)    = \lambda C \,\Delta t$
    \item $P(M_1 \to NC)    = \lambda (1-C)\,\Delta t$
    \item $P(M_1 \to F)     = \lambda (1-C)\,\Delta t$
    \item $P(M_1 \to M_1)   = 1 - 2\lambda\,\Delta t$
\end{itemize}

\subsection*{Estado $M_2$}
\begin{itemize}
    \item $P(M_2 \to M_1) = \mu_c\,\Delta t$
    \item $P(M_2 \to F)   = \lambda\,\Delta t$
    \item $P(M_2 \to M_2) = 1 - (\lambda + \mu_c)\,\Delta t$
\end{itemize}

\subsection*{Estado $FC$}
\begin{itemize}
    \item $P(FC \to M_1) = \mu_c\,\Delta t$
    \item $P(FC \to F)   = \lambda\,\Delta t$
    \item $P(FC \to FC)  = 1 - (\lambda + \mu_c)\,\Delta t$
\end{itemize}

\subsection*{Estado $NC$}
\begin{itemize}
    \item $P(NC \to M_1) = \mu_p\,\Delta t$
    \item $P(NC \to F)   = \lambda\,\Delta t$
    \item $P(NC \to NC)  = 1 - (\lambda + \mu_p)\,\Delta t$
\end{itemize}

\subsection*{Estado $F$}
\begin{itemize}
    \item $P(F \to M_1) = \mu_c\,\Delta t$
    \item $P(F \to F)   = 1 - \mu_c\,\Delta t$
\end{itemize}

% ──────────────────────────────────────────────
\newpage
\section{Derivação das Equações Diferenciais}

O passo a passo a seguir ilustra como as probabilidades de transição
definidas na seção anterior conduzem a equações diferenciais. O
procedimento é apresentado em detalhe para o estado $M_1$ e depois
generalizado.

\subsection*{Balanço de probabilidade para $M_1$ em $t + \Delta t$}

O sistema estará em $M_1$ no instante $t + \Delta t$ se, e somente se,
uma das seguintes situações ocorrer:

\begin{itemize}
    \item O sistema já estava em $M_1$ em $t$ e nenhuma falha ocorreu,
          com probabilidade $P_{M_1}(t)\cdot(1 - 2\lambda\,\Delta t)$.
    \item O sistema estava em $M_2$ em $t$ e um reparo corretivo ocorreu,
          com probabilidade $P_{M_2}(t)\cdot\mu_c\,\Delta t$.
    \item O sistema estava em $FC$ em $t$ e um reparo corretivo ocorreu,
          com probabilidade $P_{FC}(t)\cdot\mu_c\,\Delta t$.
    \item O sistema estava em $NC$ em $t$ e um reparo preventivo ocorreu,
          com probabilidade $P_{NC}(t)\cdot\mu_p\,\Delta t$.
    \item O sistema estava em $F$ em $t$ e um reparo corretivo restaurou
          o sistema, com probabilidade $P_{F}(t)\cdot\mu_c\,\Delta t$.
\end{itemize}

Somando essas parcelas:
\begin{align*}
P_{M_1}(t + \Delta t)
    &= P_{M_1}(t)\,(1 - 2\lambda\,\Delta t) \\
    &\quad + P_{M_2}(t)\,\mu_c\,\Delta t \\
    &\quad + P_{FC}(t)\,\mu_c\,\Delta t \\
    &\quad + P_{NC}(t)\,\mu_p\,\Delta t \\
    &\quad + P_{F}(t)\,\mu_c\,\Delta t.
\end{align*}

\subsection*{Passagem ao limite $\Delta t \to 0$}

Isolando o incremento de probabilidade:
\[
P_{M_1}(t + \Delta t) - P_{M_1}(t)
    = \bigl(
        -2\lambda\,P_{M_1}(t)
        + \mu_c\,P_{M_2}(t)
        + \mu_c\,P_{FC}(t)
        + \mu_p\,P_{NC}(t)
        + \mu_c\,P_F(t)
      \bigr)\Delta t.
\]

Dividindo por $\Delta t$ e tomando o limite:
\[
\frac{dP_{M_1}(t)}{dt}
    = -2\lambda\,P_{M_1}(t)
      + \mu_c\,P_{M_2}(t)
      + \mu_c\,P_{FC}(t)
      + \mu_p\,P_{NC}(t)
      + \mu_c\,P_F(t).
\]

Aplicando o mesmo procedimento aos demais estados, obtém-se o sistema
completo apresentado a seguir.

% ──────────────────────────────────────────────
\newpage
\section{Sistema de Equações Diferenciais}

A variação temporal da probabilidade de cada estado é o balanço entre
os fluxos de entrada e de saída:

\begin{alignat*}{2}
\frac{dP_{M_1}}{dt} &= -2\lambda\,P_{M_1}
    &&+ \mu_c\,P_{M_2} + \mu_c\,P_{FC} + \mu_p\,P_{NC} + \mu_c\,P_F \\[4pt]
\frac{dP_{M_2}}{dt} &= \phantom{-}\lambda C\,P_{M_1}
    &&- (\lambda + \mu_c)\,P_{M_2} \\[4pt]
\frac{dP_{FC}}{dt}  &= \phantom{-}\lambda C\,P_{M_1}
    &&- (\lambda + \mu_c)\,P_{FC} \\[4pt]
\frac{dP_{NC}}{dt}  &= \phantom{-}\lambda(1-C)\,P_{M_1}
    &&- (\lambda + \mu_p)\,P_{NC} \\[4pt]
\frac{dP_{F}}{dt}   &= \phantom{-}\lambda(1-C)\,P_{M_1}
    &&+ \lambda\,P_{M_2} + \lambda\,P_{FC} + \lambda\,P_{NC} - \mu_c\,P_F
\end{alignat*}

com a condição de normalização $\sum_i P_i(t) = 1$.

% ──────────────────────────────────────────────
\newpage
\section{Forma Matricial}

\subsection*{Notação vetorial}

Definindo o vetor linha de probabilidades
$\mathbf{P}(t) = [P_{M_1}(t),\; P_{M_2}(t),\; P_{FC}(t),\; P_{NC}(t),\; P_F(t)]$,
o sistema pode ser escrito de forma compacta como:
\[
    \dot{\mathbf{P}}(t) = \mathbf{P}(t)\,\mathbf{Q},
\]
onde $\mathbf{Q}$ é a \textbf{matriz geradora}, cujos coeficientes são as taxas de transição levantadas nas seções anteriores.

\subsection*{Matriz geradora $\mathbf{Q}$}

Na ordem de estados $\{M_1, M_2, FC, NC, F\}$:

\[
\mathbf{Q} =
\begin{bmatrix}
-2\lambda       & \lambda C          & \lambda C          & \lambda(1-C)       & \lambda(1-C) \\
\mu_c           & -(\lambda + \mu_c) & 0                  & 0                  & \lambda      \\
\mu_c           & 0                  & -(\lambda + \mu_c) & 0                  & \lambda      \\
\mu_p           & 0                  & 0                  & -(\lambda + \mu_p) & \lambda      \\
\mu_c           & 0                  & 0                  & 0                  & -\mu_c
\end{bmatrix}
\]

% ──────────────────────────────────────────────
\newpage
\section{Disponibilidade $A(t)$}

A condição inicial $\mathbf{P}(0) = \mathbf{P}_0 = [1, 0, 0, 0, 0]$ representa o sistema em plena operação no instante $t = 0$, isto é, toda a probabilidade concentrada no estado $M_1$. A disponibilidade do sistema no instante $t$ é a probabilidade de ele se encontrar em qualquer estado operacional, ou seja, em qualquer estado diferente de $F$:
\[
    A(t) = P_{M_1}(t) + P_{M_2}(t) + P_{FC}(t) + P_{NC}(t).
\]

Como as probabilidades de todos os estados somam 1 para todo $t$, essa expressão equivale a:
\[
    \boxed{A(t) = 1 - P_F(t).}
\]

Portanto, basta acompanhar a evolução de $P_F(t)$ para obter a
disponibilidade do sistema em qualquer instante.

% ──────────────────────────────────────────────
\newpage
\section{Disponibilidade em Regime Permanente}

\subsection*{Condição de equilíbrio}

Após um tempo suficientemente longo, as probabilidades $P_i(t)$ deixam de
variar. Denotando os valores estacionários por $\pi_i$, impõe-se:
\[
    \frac{dP_i}{dt} = 0, \quad \forall\, i.
\]

A condição de normalização continua válida:
\[
    \pi_{M_1} + \pi_{M_2} + \pi_{FC} + \pi_{NC} + \pi_F = 1.
\]

\subsection*{Sistema de equações de balanço}

Substituindo $dP_i/dt = 0$ no sistema da Seção~4, obtém-se:

\begin{alignat*}{2}
0 &= -2\lambda\,\pi_{M_1}
    &&+ \mu_c\,\pi_{M_2} + \mu_c\,\pi_{FC} + \mu_p\,\pi_{NC} + \mu_c\,\pi_F \\[4pt]
0 &= \phantom{-}\lambda C\,\pi_{M_1}
    &&- (\lambda + \mu_c)\,\pi_{M_2} \\[4pt]
0 &= \phantom{-}\lambda C\,\pi_{M_1}
    &&- (\lambda + \mu_c)\,\pi_{FC} \\[4pt]
0 &= \phantom{-}\lambda(1-C)\,\pi_{M_1}
    &&- (\lambda + \mu_p)\,\pi_{NC} \\[4pt]
0 &= \phantom{-}\lambda(1-C)\,\pi_{M_1}
    &&+ \lambda\,\pi_{M_2} + \lambda\,\pi_{FC} + \lambda\,\pi_{NC} - \mu_c\,\pi_F
\end{alignat*}

\subsection*{Forma matricial}

As cinco equações acima são linearmente dependentes. Por isso, substitui-se
a última delas pela condição de normalização. O sistema resultante é:

\[
\begin{bmatrix}
\pi_{M_1} \\ \pi_{M_2} \\ \pi_{FC} \\ \pi_{NC} \\ \pi_F
\end{bmatrix}^{T}
\begin{bmatrix}
-2\lambda       & \lambda C          & \lambda C          & \lambda(1-C)       & 1 \\
\mu_c           & -(\lambda + \mu_c) & 0                  & 0                  & 1 \\
\mu_c           & 0                  & -(\lambda + \mu_c) & 0                  & 1 \\
\mu_p           & 0                  & 0                  & -(\lambda + \mu_p) & 1 \\
\mu_c           & 0                  & 0                  & 0                  & 1
\end{bmatrix}
=
\begin{bmatrix}
0 \\ 0 \\ 0 \\ 0 \\ 1
\end{bmatrix}^{T}
\]

A solução desse sistema linear fornece todas as probabilidades estacionárias,
e a disponibilidade em regime permanente é:
\[
    \boxed{A_{ss} = 1 - \pi_F.}
\]

% ──────────────────────────────────────────────
\newpage
\section{Simulação por Discretização Temporal}

O vetor $\mathbf{P}^{(k)} = [P_{M_1}^{(k)},\, P_{M_2}^{(k)},\, P_{FC}^{(k)},\,
P_{NC}^{(k)},\, P_F^{(k)}]$ contém a probabilidade de o sistema estar em cada
um dos cinco estados no instante $t^{(k)} = k\,\Delta t$. No instante $k = 0$,
o sistema está em plena operação com redundância ativa, ou seja, com certeza
no estado $M_1$:
\[
    \mathbf{P}^{(0)} = [1,\; 0,\; 0,\; 0,\; 0].
\]

\subsection*{Do instante 0 para o instante 1}

Após um intervalo $\Delta t$, a probabilidade de estar em $M_1$ é a soma
das contribuições de todos os estados que podem chegar a $M_1$, ponderadas
pelas probabilidades de transição definidas na Seção~2:
\begin{align*}
P_{M_1}^{(1)} &= P_{M_1}^{(0)}\,(1 - 2\lambda\,\Delta t)
               + P_{M_2}^{(0)}\,\mu_c\,\Delta t
               + P_{FC}^{(0)}\,\mu_c\,\Delta t
               + P_{NC}^{(0)}\,\mu_p\,\Delta t
               + P_{F}^{(0)}\,\mu_c\,\Delta t.
\end{align*}

O mesmo balanço vale para cada um dos outros quatro estados. Escrevendo
os cinco balanços simultaneamente em notação matricial, obtém-se:
\[
    \mathbf{P}^{(1)} = \mathbf{P}^{(0)}\,\mathbf{T},
\]
onde $\mathbf{T}$ é a \textbf{matriz de transição} cujo elemento $(i,j)$
é a probabilidade de ir do estado $i$ para o estado $j$ em $\Delta t$.
Comparando com a matriz geradora $\mathbf{Q}$ da Seção~5 — cujos
elementos fora da diagonal são as taxas $q_{ij}$ e cuja diagonal é
$q_{ii} = -\sum_{j \neq i} q_{ij}$ — essa matriz de transição é:
\[
    \mathbf{T} = \mathbf{I} + \mathbf{Q}\,\Delta t,
\]
onde $\mathbf{I}$ é a matriz identidade $5 \times 5$. O termo
$q_{ij}\,\Delta t$ fora da diagonal é a probabilidade de transição
$i \to j$, e o termo $1 + q_{ii}\,\Delta t$ na diagonal é a
probabilidade de permanecer em $i$, conforme a Seção~2.

\subsection*{Do instante 1 para o instante 2}

Com $\mathbf{P}^{(1)}$ calculado, o mesmo balanço se repete usando
a mesma matriz $\mathbf{T}$:
\[
    \mathbf{P}^{(2)} = \mathbf{P}^{(1)}\,\mathbf{T}.
\]

\subsection*{Do instante $k$ para o instante $k+1$}

O padrão se generaliza para qualquer passo:
\[
    \mathbf{P}^{(k+1)} = \mathbf{P}^{(k)}\,\mathbf{T},
    \quad k = 0, 1, \ldots, N-1.
\]

A aproximação é válida enquanto $\Delta t$ for pequeno o suficiente para
que a probabilidade de duas transições ocorrerem no mesmo intervalo seja
desprezível, ou seja, $(\lambda\,\Delta t)(\mu_c\,\Delta t) = O(\Delta t^2)
\approx 0$ — a mesma hipótese da Seção~1.

\subsection*{Disponibilidade a cada passo}

Em cada instante $k$, a probabilidade de falha total é $P_F^{(k)}$,
a quinta componente de $\mathbf{P}^{(k)}$. A disponibilidade é:
\[
    A^{(k)} = 1 - P_F^{(k)}.
\]

À medida que $k$ cresce, $A^{(k)}$ converge para $A_{ss}$, conforme
ilustrado na Figura~\ref{fig:disponibilidade}.

\begin{figure}[h]
    \centering
    \includegraphics[width=0.85\textwidth]{disponibilidade.png}
    \caption{Disponibilidade $A(t)$ obtida pela discretização temporal com
             $\Delta t$ fixo. A linha tracejada indica o valor estacionário
             $A_{ss}$ calculado pela inversão matricial da Seção~6.}
    \label{fig:disponibilidade}
\end{figure}

\end{document}